% !TEX root = ../main.tex
\chapter{Unified Memory Model and Memory Pressure}
\label{ch:unified_memory}

\chref{ch:apple_silicon} established that Apple silicon is a
heterogeneous SoC in which every compute unit reads and writes a
single physical DRAM pool through a coherent fabric. That property has
a name --- \emph{unified memory} --- which Apple uses in marketing but
which, for a framework writer, has a precise technical meaning. This
chapter develops that meaning. The central question is not whether
memory is shared (it always is) but under what contract two different
compute units, two different command queues, two different threads
observe the same byte at the same logical instant, what the runtime is
allowed to assume about ordering and visibility, and what happens when
those assumptions interact with the macOS virtual memory subsystem.

The pay-off for \Lucid\ of getting this right is large: under unified
memory, the natural framework design has \emph{no explicit data
transfer step}. A tensor allocated by the GPU runtime can be read by
a CPU consumer, and vice versa, without a \code{cudaMemcpy}-equivalent
intermediary. \chref{ch:bridge_boundaries} catalogues the six places
where \Lucid\ touches the world outside the engine, and not one of
them is a CPU-to-GPU transfer; the bridge is internal and free. The
job of this chapter is to explain why that is correct, where its
correctness is fragile, and how \Lucid\ exposes the user-visible
consequences (\code{memory\_stats}, OOM behaviour, the choice of
device argument) in a way that lines up with the underlying
hardware.

The chapter is organised so that the reader who only wants to use
\Lucid\ may stop at \secref{sec:unified_memory:ml}; the rest is for
the reader who needs to debug a memory pathology or extend the
runtime.

\section{The Unified Memory Contract}
\label{sec:unified_memory:contract}

A \emph{memory model} is a triple \((\mathcal{A}, \mathcal{C}, \mathcal{O})\)
of address spaces, coherence guarantees, and ordering guarantees that
together define when one client of memory observes another client's
writes. The Apple silicon unified memory model is, informally:

\begin{contract}[Apple silicon unified memory]
\label{contract:unified_memory}
There is a single physical address space \(\mathcal{A}_{phys}\)
shared by every compute unit on the SoC. Each compute unit (a CPU
hardware thread, the GPU command queue, the Neural Engine, the AMX
coprocessor) sees \(\mathcal{A}_{phys}\) through a process-local
virtual address space \(\mathcal{A}_{virt}\) maintained by the macOS
kernel. Coherence is maintained at cache-line granularity by the
Apple Fabric (\secref{sec:apple_silicon:floorplan}); a write by one
client is eventually visible to every other client, and the system
provides explicit barrier primitives that turn ``eventually visible''
into ``visible by program order at this point.''
\end{contract}

The contract has three components, each of which has a
software-visible consequence for \Lucid.

\paragraph{Address spaces are unified at the physical layer.} A
buffer allocated by a CPU thread (via \code{posix\_memalign},
\code{malloc}, or the Metal allocator) lives in DRAM that the GPU
command queue can read. There is no separate GPU address space.
Apple's Metal API exposes this as the \code{MTL::Buffer} type:
buffers are allocated with a \emph{storage mode} that determines
whether the CPU has a mapped view of the same physical pages
(\secref{sec:unified_memory:storage_modes}), but no storage mode
makes the buffer \emph{inaccessible} to the GPU --- the GPU is
always the consumer.

\paragraph{Coherence is hardware, but barriers are software.} The
fabric maintains coherence at cache-line granularity, but this only
means that no two clients hold mutually incoherent copies of the same
line at the same time. It does \emph{not} mean that an arbitrary CPU
load following an arbitrary GPU command-buffer commit will observe
the GPU's writes. Visibility between the CPU and the GPU requires an
explicit barrier (\code{MTLCommandBuffer::waitUntilCompleted}, an
\code{MTLSharedEvent}, or an MLX \code{eval()} call), without which
the CPU may observe stale data even though the underlying physical
pages have been updated.

\paragraph{Ordering across clients is partial.} Within a single CPU
thread, ordering is determined by the ARM memory model (ARMv8.0-A's
\emph{weak} ordering, refined by ARMv8.3-A's release/acquire
semantics). Within a single GPU command buffer, ordering is by
program order in the encoded compute pass. Across the two, ordering
is determined exclusively by Metal synchronisation primitives. There
is no implicit happens-before relation between a Metal command
buffer's submission and a subsequent CPU load.

For \Lucid\ the practical implication is that \emph{every} CPU-side
access of a GPU-produced buffer must be preceded by either an
\code{mlx::core::eval()} on the producing array or by an explicit
Metal synchronisation primitive. The \code{SharedStorage} bridge
(\secref{sec:unified_memory:bridge}) consolidates this discipline
into a single function call so that user code does not have to
manage it. A historical bug in 3.2.1 in which \code{detach()} on a
parameter leaked an MLX lazy graph
(\chref{ch:stabilization}) was traceable to a missed \code{eval()}
on the bridge boundary; we treat that incident in
\secref{sec:unified_memory:pathologies}.

\section{Address Spaces and the macOS Virtual Memory Subsystem}
\label{sec:unified_memory:vm}

The unified physical pool is mediated, on every access, by the macOS
virtual memory subsystem. Apple silicon Macs use a page size of
\(16\) KiB (in contrast to the \(4\) KiB page size on x86 Macs). This
larger page size reduces TLB pressure --- a P-core's L1 TLB has on
the order of a few hundred entries, which translates to a few
megabytes of working set under \(4\) KiB pages but tens of megabytes
under \(16\) KiB pages --- and it matches the natural allocation
granularity of the Metal allocator.

\subsection{Page residency states}
\label{ssec:unified_memory:residency}

The macOS VM subsystem classifies every physical page into one of
several residency states, which are visible through
\code{vm\_stat(1)} and through the \code{host\_statistics64} mach
call:

\begin{itemize}
  \item \textbf{Free.} Available for new allocations.
  \item \textbf{Active.} Mapped into a process and recently
    referenced.
  \item \textbf{Inactive.} Mapped, not recently referenced, eligible
    for reclaim.
  \item \textbf{Speculative.} Prefetched but not yet referenced.
  \item \textbf{Wired.} Pinned in DRAM, cannot be paged out. The
    Metal allocator wires GPU-visible buffers because the GPU has no
    page-fault handler; an unbacked GPU access would be a hardware
    fault.
  \item \textbf{Compressed.} Belonging to a process but stored in the
    kernel's compressed memory pool rather than DRAM; reads incur
    decompression overhead.
  \item \textbf{Purgeable.} Allocated but marked volatile; the kernel
    is permitted to discard the contents under pressure.
\end{itemize}

For \Lucid, the relevant state is \emph{wired}. GPU buffers are
wired; they do not page out and do not compress. Total wired memory
is the principal hard ceiling on \Lucid's working set. On a
\(16\) GB M-series machine, sustained wired allocation above
roughly \(11\)--\(12\) GB triggers memory-pressure events that may
result in the macOS kernel terminating background processes.

\subsection{Compressed memory}
\label{ssec:unified_memory:compressed}

macOS has, since OS X Mavericks, maintained a compressed in-DRAM
pool as an alternative to swap. When the system is under pressure
and a candidate page is identified for eviction, the kernel
compresses the page (typically achieving \(2\!\!:\!\!1\) to
\(3\!\!:\!\!1\) ratios on application data) and stores it in the
compressed pool rather than writing to NVMe. A subsequent reference
decompresses the page back into uncompressed DRAM.

This is invisible to \Lucid\ at the framework level for two reasons:
\begin{enumerate}
  \item GPU-visible buffers are wired and are not eligible for
    compression.
  \item CPU-only tensors (Python integer scalars, optimiser state
    on the CPU stream, etc.) are too small for compression to be a
    significant factor.
\end{enumerate}
The compressed pool nonetheless affects sustained behaviour: when
total system pressure rises, compression overhead competes with
\Lucid's compute for the same CPU cycles, and the user observes a
slowdown without a corresponding rise in DRAM usage.

\subsection{Swap and the NVMe controller}
\label{ssec:unified_memory:swap}

Beyond compression, the kernel will swap to NVMe. On Apple silicon
this is a single-namespace NVMe SSD with very high sustained
bandwidth (multiple GB/s read, comparable write), but the latency
floor is still in the tens of microseconds. A training step that
involves a swapped page taking even one TLB miss to NVMe is several
hundred GPU kernel-launches' worth of stall. \Lucid's recommendation
to users is that training workloads be sized to fit within wired
memory plus a generous margin; the framework does nothing special to
coexist with swap, because no neural network of useful size can
afford it.

\section{Metal Resource Storage Modes}
\label{sec:unified_memory:storage_modes}

The Metal API exposes the unified memory contract through the
\code{MTLResourceStorageMode} enum, which is set at buffer creation
time and which selects how the CPU and the GPU each see the buffer's
bytes. There are four modes; \Lucid\ uses only one of them
exclusively, but understanding the others is necessary to understand
\Lucid's interoperability with the rest of the Apple graphics stack.
\tabref{tab:unified_memory:storage_modes} summarises all four for
quick reference.

\begin{table}[t]
\centering
\small
\begin{tabular}{lllll}
\toprule
Mode & CPU access & GPU access & Coherence & Lucid use \\
\midrule
\code{Shared}      & yes (mapped) & yes & fabric & \textbf{exclusive} \\
\code{Private}     & no (blit)    & yes & --     & not used \\
\code{Managed}     & yes (copy)   & yes & explicit & not used (legacy x86) \\
\code{Memoryless}  & no           & yes (tile) & --  & not used (graphics only) \\
\bottomrule
\end{tabular}
\caption[Metal storage modes.]{The four \code{MTLResourceStorageMode}
options. \Lucid\ uses only \code{Shared}, because it is the only mode
that admits a zero-copy CPU view of GPU-resident bytes.
\code{Managed} is the legacy mode for Intel Macs with discrete GPUs
and is silently promoted to \code{Shared} on Apple silicon.}
\label{tab:unified_memory:storage_modes}
\end{table}

\subsection{Shared (\code{MTLStorageModeShared})}
\label{ssec:unified_memory:shared}

In \emph{Shared} mode, the buffer's pages are mapped into both the
CPU's virtual address space and the GPU's command-queue address
space. Loads and stores from either side observe the same bytes.
Coherence is maintained by the fabric without software intervention,
modulo the visibility-by-barrier rule discussed above.

This is the storage mode that makes unified memory tangible to
software. It is also the storage mode that \MLX\ uses for nearly
every array, and consequently the mode that \Lucid's \code{Tensor}
class effectively assumes for all GPU-resident data. The
\code{SharedStorage} bridge (\secref{sec:unified_memory:bridge})
operates on \code{MTLStorageModeShared} buffers, because they are
the only buffers a CPU pointer can address.

\subsection{Private (\code{MTLStorageModePrivate})}
\label{ssec:unified_memory:private}

In \emph{Private} mode, the buffer's pages are mapped only to the
GPU. The CPU cannot read them directly; doing so requires an
explicit blit through a staging Shared buffer. Private mode exists
because the GPU may, for some configurations and some Metal device
families, place private buffers in a region of memory with different
cache attributes than shared memory --- specifically, with
write-back caching enabled on the GPU side, which is slightly faster
for purely GPU workloads.

\Lucid\ does not use \emph{Private} mode. The reasoning is precise:
the speed advantage of Private over Shared is small (single-digit
percent on Apple silicon, where the unified memory hierarchy means
both modes ultimately go to the same DRAM), and the cost is that the
zero-copy CPU bridge ceases to be zero-copy. We would have to
allocate a parallel Shared staging buffer and blit on every CPU
access. \MLX\ makes the same choice, so the question reduces to
whether \Lucid\ would deviate from \MLX, and the answer is no.

\subsection{Managed (\code{MTLStorageModeManaged})}
\label{ssec:unified_memory:managed}

\emph{Managed} mode is a legacy mode from the era of Intel Macs with
discrete GPUs (typically AMD Radeon Pro). It maintained two physical
copies of the buffer, one in VRAM and one in system DRAM, and
required explicit \code{didModifyRange} and
\code{synchronizeResource} calls to mark when each side had updated
its copy. The Apple silicon Metal driver does not support
\emph{Managed} mode on integrated GPUs (it would require maintaining
two copies in the same physical pool, which is meaningless), and
\Lucid\ never encounters it.

We mention it only because legacy \refframework\ code paths
sometimes select \emph{Managed} when the user passes
\code{device=`mps'} on macOS; this is harmless but wasteful, since
the driver silently promotes \emph{Managed} to \emph{Shared} on
Apple silicon.

\subsection{Memoryless (\code{MTLStorageModeMemoryless})}
\label{ssec:unified_memory:memoryless}

\emph{Memoryless} buffers are render-target tiles: they live only in
the GPU's on-chip tile memory and never back to DRAM at all. They are
used for intermediate render targets in deferred shading. Compute
workloads cannot use them; \Lucid\ does not encounter them.

\subsection{What \Lucid\ uses}
\label{ssec:unified_memory:lucid_mode}

\Lucid\ uses \emph{Shared} exclusively. The choice is forced by the
bridge design (\secref{sec:unified_memory:bridge}) and by \MLX's own
default. The cost is a small loss of GPU-side caching efficiency
relative to \emph{Private}; the gain is that every tensor in \Lucid\
admits a zero-copy CPU view, every interop with \Accelerate\ is
free, and the bridge contract is consistent across the entire engine.

\section{Memory Pressure on macOS}
\label{sec:unified_memory:pressure}

macOS reports system memory pressure on three discrete levels:
\begin{itemize}
  \item \textbf{Normal.} No pressure.
  \item \textbf{Warning.} The system has decided that processes
    should free non-essential caches. Reported through the
    \code{DISPATCH\_MEMORYPRESSURE\_WARN} GCD event and through the
    \code{kern.memorystatus\_vm\_pressure\_level} \code{sysctl}.
  \item \textbf{Critical.} The system is about to start terminating
    background processes to reclaim memory.
    \code{DISPATCH\_MEMORYPRESSURE\_CRITICAL}.
\end{itemize}

The kernel raises the pressure level as a function of free memory,
swap activity, and the rate of compressed-pool growth. \Lucid\ does
not subscribe to these events at the C++ engine layer, but it does
honour them indirectly: the Metal allocator, the \MLX\ buffer pool,
and \Accelerate\ all subscribe and release cached buffers when the
system signals warning or critical.

\subsection{Pressure-induced allocator behaviour}
\label{ssec:unified_memory:allocator_pressure}

\Lucid's allocator (\code{lucid/\_C/core/Allocator.\{h,cpp\}}) is a
size-class pool: requests are rounded up to one of \(23\) size
classes, and free blocks within each class are retained in a freelist
up to a configurable depth (\code{kMaxDepth = 32}). The pool is
process-local and is not coordinated with the kernel's memory
pressure events, on the principle that pressure handling is the
underlying \MLX\ and Metal allocators' responsibility; \Lucid's pool
is a layer above them and would just create a second cache layer to
manage.

This design has a consequence: \Lucid's reported memory usage
(\code{lucid.memory\_stats()}) is \emph{the pool's} reported usage,
which can be larger than the actual GPU resident set after \MLX\
has released cached buffers. The two numbers may drift by a few
hundred megabytes under sustained training; this is not a leak. We
document the distinction in \secref{sec:unified_memory:accounting}.

\subsection{User-visible pressure symptoms}
\label{ssec:unified_memory:pressure_symptoms}

For a user running \Lucid, the symptoms of memory pressure escalate
in the following sequence:

\begin{enumerate}
  \item \textbf{Slowdown without an obvious culprit.} The compressed
    pool is growing; the CPU is paying compression overhead.
  \item \textbf{Increased CPU usage in \code{kernel\_task}.} The
    kernel is spending cycles on compression / decompression.
  \item \textbf{Background apps quit silently.} Pressure has reached
    \emph{Critical} and \code{jetsam} (the OS X analog of OOM-killer)
    has begun reclaiming. \Lucid's process itself is not at risk yet
    if it is the user's foreground process.
  \item \textbf{\Lucid\ raises \code{LucidError("out of memory")}.}
    \Lucid's allocator has been refused by \MLX, which has been
    refused by Metal. At this point the only recourse is to reduce
    the working set --- smaller batch, fewer activations retained,
    gradient checkpointing.
\end{enumerate}

The framework deliberately does not silently fall back to a smaller
batch size on OOM. The reason is reproducibility: a training run
whose batch size is implicitly modulated by system memory pressure
produces different gradients on different runs, which violates the
deterministic-by-default contract documented in
\chref{ch:reproducibility}.

\section{Lucid's Memory Accounting}
\label{sec:unified_memory:accounting}

\Lucid\ exposes memory usage through a small set of Python APIs that
sit on top of the C++ engine's accounting. The public surface is
deliberately narrow:

\begin{lstlisting}[style=lucid,language=LucidPython]
import lucid

lucid.memory_stats()          # dict for current device
lucid.memory_stats("cpu")     # CPU pool stats
lucid.memory_stats("metal")   # GPU pool stats
lucid.empty_cache()           # drain the allocator freelists
\end{lstlisting}

The returned dictionary contains the following fields, all in bytes:

\begin{itemize}
  \item \code{allocated} --- total bytes currently held by live
    \code{Tensor} objects.
  \item \code{reserved} --- total bytes the allocator has from the
    underlying provider (\MLX\ or \Accelerate); always
    \(\geq\) \code{allocated}.
  \item \code{peak\_allocated} --- the high-water mark since the
    process began, or since the last \code{reset\_peak()} call.
  \item \code{peak\_reserved} --- analogous high-water mark for
    \code{reserved}.
  \item \code{num\_alloc\_calls}, \code{num\_free\_calls} ---
    operational counters for diagnostics.
\end{itemize}

\subsection{Pool versus provider}
\label{ssec:unified_memory:pool_vs_provider}

It is worth being explicit about the distinction between the
\emph{pool} (\Lucid's own freelist) and the \emph{provider}
(\MLX\ or \Accelerate). When a tensor is freed by the Python garbage
collector, its bytes return to \Lucid's pool, not to the underlying
provider. The provider sees its allocation count as constant. If a
user observes \code{lucid.memory\_stats()['allocated']} dropping to
near zero but the system's \code{Activity Monitor} continuing to
report multiple gigabytes of Metal allocation, this is the pool
holding cached buffers in case a subsequent allocation can reuse
them. \code{lucid.empty\_cache()} drains the pool back to the
provider, releasing memory to the kernel.

\subsection{Why the pool exists}
\label{ssec:unified_memory:pool_rationale}

The natural question is whether \Lucid\ needs a pool at all, given
that both \MLX\ and \Accelerate\ have their own. The answer is yes,
for two reasons. First, neither \MLX\ nor \Accelerate\ exposes the
size-class structure that \Lucid's autograd engine needs to amortise
the allocation cost of a backward pass that produces many tensors
of identical shape (one per linear layer, for example). Second, the
two providers do not coordinate with each other: a CPU-pool free
does not release memory that a GPU allocation can reuse, even though
the physical DRAM is shared. \Lucid's pool is a thin wrapper that
applies size-class amortisation uniformly and that, indirectly,
ensures that \code{empty\_cache()} drains both providers via a
single call.

\section{The Zero-Copy CPU--GPU Bridge}
\label{sec:unified_memory:bridge}

The most consequential consequence of the unified memory contract,
for \Lucid's design, is that the CPU and the GPU can share a buffer
without copying it. This section describes how that is implemented.

\figref{fig:unified_memory:bridge} sketches the three-step protocol.
The crucial observation is that no bytes move at any step --- the
\code{contiguous} call may materialise a copy if the input is
non-contiguous, but in the typical case it is a no-op, and the
remaining steps are pure metadata.

\begin{figure}[t]
\centering
\begin{adjustbox}{max width=\textwidth, center}
\begin{tikzpicture}[
  box/.style={draw, rounded corners=2pt, minimum width=3.6cm,
              minimum height=1.2cm, font=\small\sffamily,
              align=center, inner sep=4pt},
  mlx/.style={box, fill=lucid.info.bg},
  cpu/.style={box, fill=lucid.note.bg},
  buf/.style={box, fill=lucid.ok.bg, minimum width=10cm,
              minimum height=1.3cm},
  arrow/.style={->, >=stealth, thick, color=lucid.accent.dark,
                shorten >=2pt, shorten <=2pt},
  steplbl/.style={font=\footnotesize\itshape, color=lucid.muted,
               fill=white, inner sep=2pt},
]
  % MLX side
  \node[mlx] (mlx_array) at (-4, 3.5) {\texttt{mlx::core::}\\\texttt{array}};
  \node[mlx] (mlx_contig) at (-4, 1.5) {\texttt{contiguous(a)}\\\texttt{.eval()}};

  % CPU side
  \node[cpu] (cpu_view) at ( 4, 3.5) {\texttt{CpuView}\\(Lucid tensor)};
  \node[cpu] (cpu_ptr) at ( 4, 1.5) {\texttt{a.data<T>()}\\raw pointer};

  % Shared backing buffer
  \node[buf] (buffer) at (0, -0.8) {Single \texttt{MTL::Buffer} (Shared storage mode)\\physical DRAM, no copy};

  \draw[arrow] (mlx_array) -- node[steplbl, right] {step 1: ensure contiguous} (mlx_contig);
  \draw[arrow] (mlx_contig) -- node[steplbl, midway, sloped, above]
    {step 2: materialise} (buffer.north -| mlx_contig.south);
  \draw[arrow] (buffer.north -| cpu_ptr.south) -- node[steplbl, midway, sloped, above]
    {step 3: expose pointer} (cpu_ptr);
  \draw[arrow] (cpu_ptr) -- (cpu_view);
\end{tikzpicture}
\end{adjustbox}
\caption[Zero-copy bridge protocol.]{The three-step
\code{SharedStorage::from\_mlx} bridge. Step~1: ensure the MLX
array is contiguous (no-op if already, one-pass copy otherwise).
Step~2: \code{eval()} materialises the lazy graph into a single
\code{MTL::Buffer}. Step~3: \code{data<T>()} returns a CPU
pointer to that buffer. The \code{CpuView} wraps the pointer with
shape/dtype metadata, holding the MLX array alive via
\code{shared\_ptr} for the view's lifetime.}
\label{fig:unified_memory:bridge}
\end{figure}

\subsection{The bridge interface}
\label{ssec:unified_memory:bridge_iface}

The bridge is implemented in the \code{SharedStorage} class
(\code{lucid/\_C/backend/gpu/SharedStorage.\{h,cpp\}}). Its public
interface is, schematically:

\begin{lstlisting}[style=lucidcpp,language=C++]
class SharedStorage {
public:
  // Adopt an MLX array; return a CPU-readable pointer to its bytes.
  static CpuView from_mlx(const mlx::core::array& a);

  // Adopt a CPU buffer; return an MLX array that wraps the same bytes.
  static mlx::core::array from_cpu(void* data, Shape shape, Dtype dt);

  // Synchronisation point: ensure all GPU writes to `a` are visible.
  static void eval_into_cpu(const mlx::core::array& a);
};
\end{lstlisting}

The crucial operation is \code{from\_mlx}. Given an
\code{mlx::core::array}, it returns a view of the same physical
bytes that the CPU side of \Lucid\ can dereference. There is no
copy. The implementation has three steps:

\begin{enumerate}
  \item Call \code{a.eval()} (transitively, \code{mlx::core::eval})
    to materialise the array. \MLX\ is a lazy graph runtime; an
    unmaterialised array has no backing buffer.
  \item Call \code{a.data<T>()} to obtain a raw pointer to the
    backing \code{MTL::Buffer}'s contents (the buffer is
    \emph{Shared} storage mode, so the pointer is CPU-dereferenceable
    immediately).
  \item Wrap the pointer in a \code{CpuView} carrying the shape and
    dtype metadata, but \emph{ignoring} any non-contiguous strides
    from \MLX's lazy view. This is the subtle point.
\end{enumerate}

\subsection{The stride trap}
\label{ssec:unified_memory:stride_trap}

\MLX\ supports lazy view operations: \code{transpose}, \code{slice},
and \code{reshape} return arrays that share storage with their
inputs but carry adjusted strides. The bytes underlying such an
array are in input-original order, not the post-view order.

Crucially, \code{mlx::core::array::data<T>()} returns a pointer to
the raw bytes \emph{without} applying the view's strides. A naive
bridge that handed those bytes to a CPU consumer expecting
post-transpose order would corrupt every downstream computation.

The discipline \Lucid\ enforces is therefore:

\begin{lucidwarn}
Before \code{SharedStorage::from\_mlx}, the array must be made
contiguous. The bridge calls \code{mlx::core::contiguous(a).eval()}
on its way in. The order matters: \code{contiguous} first (so that
strides are unit-stride), then \code{eval} (so that the contiguous
copy is materialised), then \code{data<T>()}.
\end{lucidwarn}

This discipline is encoded in \code{SharedStorage::from\_mlx} and is
not negotiable. A historical engine quirk
(\code{engine-mlx-data-ignores-strides} in the vault) is the
source-of-truth note that captured the bug class.

\subsection{Why this is free}
\label{ssec:unified_memory:bridge_free}

The bridge is sometimes free, sometimes not, depending on whether
the source array is already contiguous:

\begin{itemize}
  \item If the array is contiguous (no transpose, no non-trivial
    slicing in its history), \code{contiguous} is a no-op: \MLX\
    recognises the trivial case and returns the same backing
    storage. The bridge is free.
  \item If the array is non-contiguous, \code{contiguous} allocates
    a new buffer and performs a kernel-driven copy. The bridge is
    not free: it costs one full read--write pass over the tensor.
    This is, however, fundamental --- no possible bridge could be
    free here, because the CPU consumer needs unit-stride data.
\end{itemize}

\Lucid's design ensures that most tensors at the bridge boundary are
already contiguous. The dispatch layer (\chref{ch:backend_dispatch})
prefers to keep operations contiguous; views are localised; the
bridge is invoked at function boundaries (e.g.\ \code{.numpy()},
\code{.tolist()}, \code{\_\_dlpack\_\_}) where the user has signalled
intent to interoperate with external code that requires unit-stride.

\section{Memory Pressure and ML Workloads}
\label{sec:unified_memory:ml}

We now turn to the practical implications for training workloads.

\subsection{The working-set decomposition}
\label{ssec:unified_memory:working_set}

A training step's working set decomposes into the following classes:

\begin{enumerate}
  \item \textbf{Model parameters.} For a parameter count \(N\) at
    FP32 precision, this is \(4N\) bytes. A 100M-parameter model
    occupies 400 MB.
  \item \textbf{Optimiser state.} For Adam, two additional copies of
    each parameter (the first and second moments), so \(8N\) bytes.
  \item \textbf{Gradients.} One copy per parameter during the
    backward pass, so an additional \(4N\) bytes peak.
  \item \textbf{Activations.} Highly model- and batch-dependent. For
    a Transformer of hidden size \(d\), depth \(L\), batch size
    \(B\), sequence length \(T\), the activation memory is roughly
    \(O(B L T d)\) for feed-forward and \(O(B L T^2)\) for
    attention. Activations are the dominant term for non-trivial
    sequence lengths.
  \item \textbf{Allocator headroom.} \(\sim 10\!-\!20\,\%\) overhead
    held by \Lucid's pool and \MLX's pool for amortisation.
\end{enumerate}

For a 100M-parameter Transformer at batch size 32 and sequence
length 512, the breakdown is approximately:
\begin{itemize}
  \item Parameters: 400 MB.
  \item Adam state: 800 MB.
  \item Gradients (peak): 400 MB.
  \item Activations: 2--4 GB depending on FFN size.
  \item Allocator headroom: 0.4--0.8 GB.
\end{itemize}
The total is in the 4--6 GB range, comfortable on a 16 GB unified
memory configuration but tight on 8 GB. Doubling the batch size
roughly doubles the activation term and pushes the total above 8 GB.

\subsection{Activation memory and gradient checkpointing}
\label{ssec:unified_memory:checkpointing}

The activation term is the largest and the most adjustable. For each
layer, the forward pass produces an activation tensor that the
backward pass consumes. Naively, all activations are retained until
backward; this gives the \(O(BLTd)\) figure above.

Gradient checkpointing (\chref{ch:gradient_checkpointing}) trades
compute for memory: by retaining only every \(k\)-th activation and
recomputing the intermediate ones on the backward pass, the
activation memory becomes \(O(BLTd / k)\) at the cost of an extra
\(k\)-fraction of forward work. \Lucid\ exposes this through
\code{lucid.utils.checkpoint}.

The break-even point for checkpointing on Apple silicon is
characteristically different from a discrete-GPU system because the
GPU's compute is relatively scarce and bandwidth is relatively
abundant. The same model that benefits from \(k = 4\) checkpointing
on a discrete GPU may prefer \(k = 2\) on M-series, because the
re-forward pass is more expensive in relative terms.

\subsection{The non-preallocation policy}
\label{ssec:unified_memory:no_prealloc}

\Lucid\ deliberately does \emph{not} preallocate a fixed memory
arena at startup. \refframework\ on CUDA does (the
\code{PYTORCH\_CUDA\_ALLOC\_CONF} environment variable can request
it), partly to avoid CUDA's fragmentation behaviour and partly to
reduce allocator overhead. On Apple silicon the fragmentation
argument is weaker (the unified pool is shared with everything else
on the machine, so a giant preallocation would stall other system
work), and the allocator overhead argument is also weaker (the size
classes in \Lucid's pool already amortise away most of the
allocation cost).

The result is that \Lucid's memory usage scales with the working
set, not with a fixed budget. This makes it easier to coexist with
other applications on the same machine (the user's IDE, their
browser, a notebook viewer) without manual tuning.

\section{Pathologies and Failure Modes}
\label{sec:unified_memory:pathologies}

The unified memory contract is strong, but two classes of bugs have
recurred in \Lucid's history. We document both because the same
patterns are likely to recur in future extensions.

\subsection{The lazy-graph leak (3.2.1)}
\label{ssec:unified_memory:lazy_leak}

The 3.2.1 series shipped a hotfix for a long-running memory leak in
the BatchNorm parameter-update path. The symptom was that, on a
ResNet-18 training loop, total Metal allocation grew by approximately
37 MB per iteration without bound, eventually causing OOM.

The root cause was an interaction between \MLX's lazy graph and
\Lucid's \code{detach()} semantics. \code{detach()} on a parameter
clears the autograd metadata (\code{grad\_fn}, \code{requires\_grad})
but does not, by itself, force the underlying \MLX\ array to
materialise. If a subsequent operation reads the detached tensor and
the operation is itself lazy, the new graph node retains a reference
to the old graph node. Across many iterations the lazy graph chain
grows without bound and the underlying buffers cannot be freed.

The fix was to call \code{eval\_tensors} on the BatchNorm running
statistics at the end of each forward pass in training mode. This
forces materialisation, breaks the lazy graph chain, and allows the
allocator to release intermediate buffers. The discipline has since
been generalised: any module that maintains an in-place running
statistic (BatchNorm, RunningMean, EMA optimiser state) is
responsible for evaluating that statistic at a known synchronisation
point.

\subsection{The InstanceNorm OOM (3.2.2)}
\label{ssec:unified_memory:instancenorm_oom}

A separate incident in 3.2.2 manifested as an OOM during
InstanceNorm forward on inputs larger than approximately
\((B, C, H, W) = (32, 64, 512, 512)\). The cause was that
InstanceNorm's forward kernel was materialising several
broadcast-expanded intermediates --- per-channel mean, per-channel
variance, broadcast-back to \((B, C, H, W)\) for the normalisation
step --- that each held a separate buffer.

The fix was algorithmic: rewrite the forward to fold the broadcast
into the normalisation kernel directly, eliminating two of the four
intermediate buffers. The result was a 4$\times$ reduction in peak
memory at the cost of a slightly more complex kernel. We treat this
case in detail in \chref{ch:no_compile_sweep}, since it is one of
several similar rewrites that improved \Lucid's normalisation family.

The general lesson is that, under unified memory, peak memory is
often dominated by short-lived intermediates rather than long-lived
state. Optimising peak therefore involves \emph{kernel} restructuring
rather than \emph{allocator} tuning.

\subsection{Cross-process sharing is not supported}
\label{ssec:unified_memory:no_cross_process}

A subtler pathology, occasionally raised by users: can two \Lucid\
processes share a tensor through \code{multiprocessing.shared\_memory}
or via the macOS \code{IOSurface} API? The answer is no, on two
grounds:

\begin{enumerate}
  \item \MLX's allocator does not produce \code{MTLBuffer}s with
    cross-process sharing flags set. The buffer's underlying pages
    are accessible only to the owning process.
  \item Even if cross-process \code{IOSurface}-backed buffers were
    created manually, \Lucid's autograd metadata is per-process; a
    tensor passed to another process would lose its \code{grad\_fn}.
\end{enumerate}

Multi-process training on a single machine would therefore require
explicit serialisation (which costs the bridge cost discussed above)
and is not currently a supported mode. Distributed training is out
of scope for \Lucid\ 3.x and is left to a future release.

\section{Summary and Forward References}
\label{sec:unified_memory:summary}

Unified memory on Apple silicon is a precisely specified contract:
single physical pool, fabric coherence, software-visible barriers,
no implicit CPU/GPU ordering. \Lucid\ exploits this contract through
the \code{SharedStorage} bridge to eliminate explicit data movement
between the CPU and GPU streams, modulo a single contiguity
requirement that the bridge enforces internally. The macOS virtual
memory subsystem adds residency states, compressed memory, and a
three-level pressure signal that \Lucid\ honours indirectly through
the providers it depends on.

The chapters that follow build on this foundation.
\chref{ch:mlx_runtime} treats \MLX\ as the GPU-side runtime that
holds most \Lucid\ buffers in \emph{Shared} mode and exposes them
through \code{array::data<T>()}. \chref{ch:accelerate} treats
\Accelerate\ as the CPU-side counterpart that consumes the same
buffers through unit-stride CPU pointers. \chref{ch:tensor_model}
formalises \Lucid's \code{Tensor} and \code{Storage} types and shows
how the bridge slots into the tensor model. Finally,
\chref{ch:bridge_boundaries} enumerates the six places in \Lucid\
where unified memory contacts the world outside the framework, and
shows why those six are the only places that ever need to perform an
explicit copy.

\begin{lucidnote}
The two operational levers a user has are
\code{lucid.empty\_cache()} (drain the allocator pool back to the
provider) and the choice of \code{device=`cpu'} vs.\
\code{device=`metal'} on tensor creation. The first lever is
non-destructive and useful when memory pressure rises. The second
lever determines which provider's pool the tensor occupies, but
\emph{not} whether the bytes themselves are shared --- under
unified memory, both providers see the same DRAM.
\end{lucidnote}
